\section{Research Gaps and Directions}
\label{sec:discussion}
\paragraph{Measuring non-target preservation.}
Image-quality metrics are insensitive to many design-specific failures: a corrected headline can shift a logo, modify a brand color, or alter the line breaks of an unrelated paragraph.
Editing benchmarks should include protected-object masks, semantic content checks, and explicit layout tolerances.
Preservation must be evaluated on the correct level: pixels for raster methods, paths for \ac{svg}, object identities for layered documents, and \ac{dom}/content structure for code-native artifacts.
However, some non-target movement is necessary: \ac{infoedit} demonstrates that infographic editing may require coordinated global reflow \citep{yang2026infoedit}.
An evaluation protocol should distinguish permitted dependent changes from changes to explicitly protected content (Eq.~\ref{eq:reflow}).

\paragraph{Representation-dependent evidence gaps.}
A generated layout may satisfy geometric constraints without carrying the text and image assets needed for document revision.
A separated raster layer stack may support recomposition while losing semantic text or stable design-object identities, as illustrated by raster-to-layer decomposition \citep{layerd2025}.
Conversely, vector paths or code-native objects may remain manipulable even when reported evaluation emphasizes appearance or document validity rather than edit-execution accuracy \citep{vectorfusion2023,designascode2026}.
The relevant question is therefore not whether a system is simply ``editable'', but which operations its delivered representation permits, which operations the authors demonstrate, and which of those operations receive quantitative evaluation.
An image-space editor may satisfy a controlled edit instruction without offering native document export; the image-input baselines in \ac{mildedit} highlight why this format limitation should not be confused with failure at the evaluated raster editing task \citep{mildedit2026}.

\paragraph{Native artifact validity.}
An image preview cannot certify that a \ac{pdf}, \ac{svg}, \ac{pptx}, \ac{psd}, or \ac{html}/\ac{css} design remains editable.
Reliable evaluation should include a round-trip: reopen the saved result in a target editor or renderer, inspect whether objects retain the expected type, perform a controlled edit, then render again.
This is especially important for code-native agents that may output visually convincing pages with brittle styles or missing assets \citep{designascode2026,editablevisual2026}.

\paragraph{Evidence-aligned comparisons.}
Cross-system comparisons must start with the task, permitted inputs, delivered artifact, and operations actually measured.
Penalizing raster editing for not exporting vector nodes, or rewarding any \ac{svg} merely because it contains paths, confuses format requirements with edit-task success.
Every editing-capability claim should therefore be linked to an output format, a concrete operation, and the relevant workflow demonstration or measurement.
For example, alpha-mask quality supports compositing readiness but not semantic text replacement; successful SVG rendering establishes visual validity but not object-identity preservation after a later edit.
We recommend reporting instruction compliance, protected-content preservation, artifact validity, visual quality, and editing effort separately, with inapplicable dimensions explicitly identified rather than silently scored as failures.

\paragraph{A minimal cross-representation editing test.}
A useful paired evaluation would start from the same rendered design and specify the same user intent---for example, replace a headline while preserving a logo and an unrelated caption---but allow each system to operate on its declared input and output representation.
The test must record whether a method receives only an image, a layer stack, or the original source document, since these inputs expose different information.
Each submission should return both a rendered result and, where its output contract promises one, a reopenable native artifact.
Scoring should distinguish exact text correctness, permissible downstream reflow, changes to explicitly protected objects, and format-validity failures.
Results should be reported by representation and intervention type before any aggregate is attempted; an unavailable native-file dimension should be marked inapplicable rather than silently scored as an editing failure.
This proposal defines a future comparison protocol, not an experiment conducted in this survey.

\paragraph{Human-centred outcomes.}
Designer workflows involve more than one-shot aesthetic ratings.
\ac{vinci} provides an example of measuring editing efficiency \citep{vinci2021}.
Future comparisons should additionally test task completion time, correction burden, how easily users can localize an error, compatibility with editing software, and designer trust across multiple changes.
User expertise and task familiarity should be treated as experimental variables.

\paragraph{Reliable evidence synthesis.}
Evaluation descriptions found in abstracts should not automatically be mapped to specific numerical scores.
Full experimental tables and appendix details can reveal sample-size limits, unreported exclusions, or test conditions that change the interpretation.
The remaining survey work therefore includes systematically logged searches and source audits, explicit boundary adjudication, bibliographic reconciliation, and separate flags for measured versus merely claimed editing capability.

\paragraph{Reproducibility and provenance.}
Source code, weights, evaluation data, and proprietary \ac{api} dependencies introduce independent reproducibility constraints.
A system may publicly release an inference pipeline without the training recipe or may require a paid closed model at evaluation time.
These attributes should be tracked separately from whether a final artifact is editable, and should not be inferred merely from an available project-page link.
